\Needspace{8\baselineskip}\originalchapter{习题课: 相对论运动学与几何截面}
本章收录Recitation 1、2及Recitation 3题2。它们训练同一项判断：先选守恒的四动量不变量，再区分束流总能量、动能和质心能量。除显式写$c$的式子外取$c=1$。每题保留原课编号；据Tianyu Justin Yang署名官方答案重编的部分明确标注，源解答不足或有误处另行补解。
\originalsection{质量、反冲与阈值}
\begin{exercise}\label{ex:rec-1-1}
Recitation 1题1。CERN的LEP将电子和正电子各加速到100GeV，求Lorentz因子$\gamma$。
\end{exercise}
\begin{solution}
官方解答重编。总能量$E=\gamma m_ec^2$，故$\gamma=10^5/0.510999\simeq1.9570\times10^5$；官方用$0.5$MeV得到$2\times10^5$。这里100GeV是单束能量，不能再把两束之和代入单个电子的$\gamma$。
\end{solution}
\begin{exercise}\label{ex:rec-1-2}
Recitation 1题2。把氘核拆成自由质子和中子，至少需要多少能量？
\end{exercise}
\begin{solution}
官方质量差方法重编并改正符号。需要补入的结合能为$B_d=(m_p+m_n-m_d)c^2\simeq2.2246$MeV，正号来自自由两核子的质量和较大。官方答案列$2.01355,1.00728,1.00867$u并报2.23MeV，但差式写反，且这些截断数直接相减为$0.00240$u，不能支持其所写$0.0025$u。应先用足够精度的质量再取有效数字。此处问分离阈值；辐射俘获的光子还承担动量，反冲修正见习题\ref{ex:8701-6-3}，两种能量不要混同。
\end{solution}
\begin{exercise}\label{ex:rec-1-3}
Recitation 1题3。一个激发粒子发射光子，该光子在什么条件下能够再次被吸收？
\end{exercise}
\begin{solution}
官方发射反冲推导重编，吸收条件为编者补全。设激发态、基态静质量分别$M,m$，初始激发粒子静止。末态动量大小均为$E_\gamma/c$，在壳条件给
\begin{equation}
 E_{\rm em}=\frac{(M^2-m^2)c^2}{2M}
 =\Delta\left(1-\frac{\Delta}{2Mc^2}\right),\qquad \Delta=(M-m)c^2.
\end{equation}
若受体是静止、相同的自由基态粒子，$p+ k$必须落在激发态质量壳上，故
\begin{equation}
 E_{\rm abs}=\frac{(M^2-m^2)c^2}{2m}
 =\Delta\left(1+\frac{\Delta}{2mc^2}\right)>\Delta>E_{\rm em}.
\end{equation}
两端都有反冲损失，所以并非“能级差相同就必然共振”。受体运动时应将光子变换到其静止系，要求$\gamma E_{\rm em}(1-\beta\cos\vartheta)=E_{\rm abs}$，在谱线宽度内匹配；也可由环境承担反冲，例如足够大质量的晶格使反冲移位小于线宽。若有其他吸收能级，其能量与动量条件也必须同时满足。
\end{solution}
\begin{exercise}\label{ex:rec-1-4}
Recitation 1题4。质子束打静止质子，产生反质子的最小束流能量是多少？
\end{exercise}
\begin{solution}
官方阈值结果重编。重子数初态为2，最低质量末态必须是$ppp\bar p$。固定靶$s=2m_p^2+2m_pE_b$，阈值所有末粒子在质心系无相对运动，$s_{\rm th}=(4m_p)^2$。所以$E_{b,\rm th}=7m_p\simeq6.568$GeV，而动能$T_{b,\rm th}=6m_p\simeq5.630$GeV。多出的束流能量用于实验室中整体末态的运动；不能只拿$2m_p$的新生对静能作阈值。
\end{solution}
\begin{exercise}\label{ex:rec-1-5}
Recitation 1题5。假设静止pion衰变为电子和正电子，两产物的速度是多少？
\end{exercise}
\begin{solution}
官方解答重编。电荷守恒选择中性$\pi^0$，这是题设的一条稀有衰变而非其主要$\gamma\gamma$道。两末态反向等动量，各有$E_e=m_{\pi^0}/2\simeq67.49$MeV，故
$v/c=\sqrt{1-4m_e^2/m_{\pi^0}^2}\simeq0.9999713$，$\gamma_e\simeq132.07$。
\end{solution}
\begin{exercise}\label{ex:rec-1-6}
Recitation 1题6。质子束与静止质子碰撞，产生$p+n+\pi^+$的最小束流能量是多少？
\end{exercise}
\begin{solution}
官方方法重编并保留精确质量阈值。
$s_{\rm th}=(m_p+m_n+m_{\pi^+})^2$，因此
\begin{equation}
 E_{b,\rm th}=\frac{(m_p+m_n+m_{\pi^+})^2-2m_p^2}{2m_p}
 \simeq1.23057\mathrm{GeV},\quad T_{b,\rm th}\simeq292.30\mathrm{MeV}.
\end{equation}
官方粗略质量与速度四舍五入报1.24GeV，与本式同一物理阈值。若说“束流能量”未声明总能量或动能，应把两者一起给出。
\end{solution}
\begin{exercise}\label{ex:rec-1-7}
Recitation 1题7。光子能量$E=h\nu=hc/\lambda$。求Compton散射中波长的变化。原题第二个等号漏$c$；以下受体为初始静止的自由电子，散射角为$\theta$。
\end{exercise}
\begin{center}\begin{tikzpicture}[>=Stealth,scale=.8]
\draw[->,photon] (-2,0)--(0,0) node[midway,above] {$k,\lambda$};
\draw[->,photon] (0,0)--(2,1.2) node[right] {$k',\lambda'$};
\draw[->,thick] (0,0)--(1.3,-1) node[right] {$p'$};
\fill (0,0) circle (1.5pt);\draw (1,0) arc (0:31:1);
\node at (1.25,.25) {$\theta$};\node[below left] at (0,0) {$p=(m_ec,\mathbf0)$};
\end{tikzpicture}\end{center}
\begin{solution}
官方四动量方法重编。设$p+k=p'+k'$，对$p'=p+k-k'$平方，利用$p^2=p'^2=m_e^2c^2$、$k^2=k'^2=0$，得到
$m_ec^2(E-E')=EE'(1-\cos\theta)$。除以$EE'$并代$E=hc/\lambda$，
\begin{equation}\lambda'-\lambda=\frac{h}{m_ec}(1-\cos\theta).\end{equation}
前向散射不变，反向最大移位$2h/(m_ec)$；这两个极限同时检查符号和因子2。
\end{solution}
\originalsection{宇宙线与散射不变量}
\begin{exercise}\label{ex:rec-2-1}
Recitation 2题1。GZK机制认为高能宇宙线质子会通过$p+\gamma_{\rm CMB}\to p+\pi^0$损失能量，背景光子$E_{\rm CMB}\simeq3\times10^{-4}$eV。取质子能量远大于质量、与光子正面相撞，估算反应阈值并与原件PDG宇宙线谱图29.8比较。
\end{exercise}
\begin{solution}
官方扫描手写不变量推导重编；传播模型限制由编者补充。初态
$s=m_p^2+2\varepsilon(E_p+p_p)$，正面碰撞及$p_p\simeq E_p$给$s\simeq m_p^2+4\varepsilon E_p$。阈值$s=(m_p+m_{\pi^0})^2$，故
\begin{equation}
 E_{p,\rm th}\simeq\frac{2m_pm_{\pi^0}+m_{\pi^0}^2}{4\varepsilon}
 \simeq2.26\times10^{20}\mathrm{eV}.
\end{equation}
这里MeV质量乘$10^6$后再与eV光子相配。原图在几$10^{19}$eV之后显示通量压低，而不是在这个单光子阈值处形成锐利、绝对的能量上限。CMB有能量分布和高能尾部，宇宙线经历多次作用与有限传播距离；要算衰减长度，须再积分光谱、夹角与能量依赖截面。不能把一次头碰近似替代传播模型，近源粒子也不受“绝无超过阈值事件”的限制。原PDG图标有CC排除条款，本稿保留比较结果与链接，不转载该图。
\end{solution}
\begin{exercise}\label{ex:rec-2-2}
Recitation 2题2。反应$1+2\to3+4$的四动量为$p_i$。讨论质心系与粒子2静止的固定靶系中的$s,t,u$，采用相对散射能量可忽略的粒子质量近似。
\end{exercise}
\begin{solution}
官方答案重编并限定近似。定义$s=(p_1+p_2)^2$、$t=(p_1-p_3)^2$、$u=(p_1-p_4)^2$。先保持在壳质量可严格证明
$s+t+u=\sum_i m_i^2$，证明只需展开平方并代$p_1+p_2=p_3+p_4$。
质心系近似各能量$E_*=\sqrt s/2$，若$\theta_*$为1到3的夹角，
\begin{equation}
 s=4E_*^2,\quad t=-\frac s2(1-\cos\theta_*),\quad
 u=-\frac s2(1+\cos\theta_*).
\end{equation}
固定靶中不能把静止受体质量设为零：$s=m_1^2+m_2^2+2m_2E_1\simeq2m_2E_1$；由$t=(p_4-p_2)^2$、$u=(p_3-p_2)^2$得
$t=m_4^2+m_2^2-2m_2E_4\simeq2m_2(E_3-E_1)$，
$u=m_3^2+m_2^2-2m_2E_3\simeq-2m_2E_3$。
固定$s$下$t$范围为$-s\le t\le0$，而非官方文字不加固定$s$限定的$-\infty<t\le0$。
\end{solution}
\begin{exercise}\label{ex:rec-3-2}
Recitation 3题2。质量$m$、半径$r$的小球束与大球作弹性硬球散射，大球质量$M\gg m$、半径$R\gg r$。
(a)求微分及总几何截面；(b)求散射角$\theta\le30^\circ$的比例；
(c)HERA电子--质子总截面约10mb，与质子半径0.862fm的几何截面比较并解释。原题$\sqrt s$括号为空，不另造束能条件。
\end{exercise}
\begin{center}\begin{tikzpicture}[>=Stealth,scale=.8]
\def\hardballradius{1.2}
\draw (0,0) circle (\hardballradius);\draw[dashed] (-2,.7)--(-.97,.7);
\draw[->,thick] (-2,.7)--(-.97,.7) node[midway,above] {入射};
\draw[->,thick] (-.97,.7)--(-1.6,1.8) node[above] {反射};
\draw[dashed] (-.97,.7)--(-.97,0);\draw[<->] (-2,0)--(-2,.7) node[midway,left] {$b$};
\draw[->] (0,0)--(-.97,.7) node[midway,below] {$a=R+r$};
\draw[dotted] (-.97,.7)--(.7,.7);\node at (-1.45,1.03) {$\theta$};
\end{tikzpicture}\end{center}
\begin{solution}
官方解答重编，数值按原题半径复算。取有效接触半径$a=R+r$；镜面反射几何给$b=a\cos(\theta/2)$，环带面积$\dd\sigma=2\pi b|\dd b|$，对应$\dd\Omega=2\pi\sin\theta\dd\theta$，因此
\begin{equation}
 \frac{\dd\sigma}{\dd\Omega}=\frac b{\sin\theta}\left|\frac{\dd b}{\dd\theta}\right|
 =\frac{a^2}{4},\qquad \sigma_{\rm tot}=\pi a^2\simeq\pi R^2.
\end{equation}
小角比例为$[1-\cos(\pi/6)]/2=0.06699$。
$\pi(0.862\mathrm{fm})^2=2.334\mathrm{fm}^2=23.34$mb，所以比值约0.428。官方文字一处写0.88fm，却使用2.33fm$^2$，本稿按题面0.862fm统一。质子电荷半径不是吸收硬球半径；落在电荷分布范围内的电子也不必发生被实验计入的散射，量子振幅和探测定义决定截面。
\end{solution}
